% =============================================================================
%  Soutenance MODAL — Modélisation stochastique d'un carnet d'ordres
%  Ettahiri Anass & Ramlaoui Mohamed Taha — 25 minutes
% =============================================================================
\documentclass[aspectratio=169,10pt]{beamer}

\usepackage[french,noconfigs]{babel}
\usepackage[utf8]{inputenc}
\usepackage[T1]{fontenc}
\usepackage{amsmath,amssymb,mathtools}
\usepackage{graphicx}
\usepackage{booktabs,array}
\usepackage{xcolor}
\usepackage{tikz}

% ── Couleurs Polytechnique ────────────────────────────────────────────────────
\definecolor{bleu303}{RGB}{0,58,86}
\definecolor{bleu315}{RGB}{0,106,140}
\definecolor{violet2622}{RGB}{89,15,85}
\definecolor{vert575}{RGB}{88,116,12}
\definecolor{marron871}{RGB}{174,154,99}
\definecolor{rouge485}{RGB}{228,32,14}
\definecolor{gris10}{RGB}{134,133,134}
\definecolor{gris428}{RGB}{218,220,222}

% ── Thème Beamer minimal (sans beamerx) ──────────────────────────────────────
\usetheme{default}
\usecolortheme[named=bleu303]{structure}
\setbeamertemplate{navigation symbols}{}
\setbeamertemplate{itemize items}{\textbullet}
\setbeamertemplate{enumerate items}[default]

% Titre de frame
\setbeamertemplate{frametitle}{%
  \vspace*{8pt}%
  \noindent\textcolor{bleu303}{\sffamily\large\bfseries\MakeUppercase{\insertframetitle}}%
  \vspace*{2pt}%
  \newline\textcolor{marron871}{\rule{\textwidth}{0.6pt}}%
  \vspace*{4pt}%
}

% Pied de page
\setbeamertemplate{footline}{%
  \hspace*{10pt}%
  \raisebox{3pt}{\textcolor{gris10}{\tiny\sffamily
    ÉCOLE POLYTECHNIQUE \quad—\quad
    Carnet d'ordres SNA}}%
  \hfill%
  \raisebox{3pt}{\textcolor{gris10}{\tiny\sffamily\insertframenumber}}%
  \hspace*{10pt}%
  \vspace*{3pt}%
}

% Blocs compacts
\setbeamertemplate{block begin}{%
  \vspace{5pt}%
  \begin{beamercolorbox}[wd=\textwidth,sep=5pt]{block body}%
  \small%
}
\setbeamertemplate{block end}{\end{beamercolorbox}\vspace{5pt}}
\setbeamercolor{block body}{bg=bleu303!8,fg=black}

% Retrait des marges latérales
\setbeamersize{text margin left=14pt, text margin right=14pt}

% ── Commande figure ───────────────────────────────────────────────────────────
\newcommand{\fig}[2][width=\linewidth]{%
  \includegraphics[#1,keepaspectratio]{../final_report/figures/#2}%
}

% ── Raccourcis maths ─────────────────────────────────────────────────────────
\newcommand{\E}{\mathbb{E}}
\newcommand{\Pb}{\mathbb{P}}
\newcommand{\IG}{\mathrm{IG}}

% ── Séparateur de partie ─────────────────────────────────────────────────────
\newcommand{\partframe}[3]{%  #1=couleur fond, #2=label (PARTIE N), #3=titre
  \setbeamertemplate{footline}{}%
  \begin{frame}[plain,c,noframenumbering]%
    \begin{tikzpicture}[overlay,remember picture]
      \fill[#1] (current page.south west) rectangle (current page.north east);
      \fill[marron871]
        (current page.north west) rectangle
        ([xshift=5pt]current page.south west);
      \node[white,font=\fontsize{24}{30}\selectfont\sffamily\bfseries,
            text width=18cm, align=center]
        at (current page.center) {#3};
      \draw[marron871!60!white,line width=0.7pt]
        ([xshift=-3.5cm,yshift=-1.2cm]current page.center)
        -- ([xshift=3.5cm,yshift=-1.2cm]current page.center);
      \node[white,opacity=0.65,font=\fontsize{9}{9}\selectfont\sffamily]
        at ([yshift=-1.9cm]current page.center)
        {\MakeUppercase{#2}};
    \end{tikzpicture}
  \end{frame}%
  \setbeamertemplate{footline}{%
    \hspace*{10pt}%
    \raisebox{3pt}{\textcolor{gris10}{\tiny\sffamily
      ÉCOLE POLYTECHNIQUE \quad—\quad Carnet d'ordres SNA}}%
    \hfill%
    \raisebox{3pt}{\textcolor{gris10}{\tiny\sffamily\insertframenumber}}%
    \hspace*{10pt}%
    \vspace*{3pt}%
  }%
}

% ── Titrage ───────────────────────────────────────────────────────────────────
\title{Modélisation stochastique d'un carnet d'ordres}
\author{Ettahiri Anass \& Ramlaoui Mohamed Taha}
\date{Mai 2026 \quad — \quad MODAL}

% =============================================================================
\begin{document}
% =============================================================================

% ── Titre ─────────────────────────────────────────────────────────────────────
{
\setbeamertemplate{footline}{}
\begin{frame}[plain,c,noframenumbering]
  \begin{tikzpicture}[overlay,remember picture]
    \fill[bleu303] (current page.south west) rectangle (current page.north east);
    \fill[marron871]
      (current page.north west) rectangle
      ([xshift=5pt]current page.south west);
    \node[white,font=\fontsize{22}{28}\selectfont\sffamily\bfseries,
          text width=16cm,align=center]
      at ([yshift=0.8cm]current page.center)
      {MODÉLISATION STOCHASTIQUE D'UN CARNET D'ORDRES};
    \draw[marron871!70!white,line width=0.7pt]
      ([xshift=-4cm,yshift=-0.5cm]current page.center)
      -- ([xshift=4cm,yshift=-0.5cm]current page.center);
    \node[white,opacity=0.75,font=\fontsize{10}{14}\selectfont\sffamily]
      at ([yshift=-1.1cm]current page.center)
      {Ettahiri Anass \& Ramlaoui Mohamed Taha};
    \node[white,opacity=0.55,font=\fontsize{9}{9}\selectfont\sffamily]
      at ([yshift=-1.9cm]current page.center)
      {MODAL — Mai 2026};
  \end{tikzpicture}
\end{frame}
}

% ── Plan ─────────────────────────────────────────────────────────────────────
\begin{frame}{Plan}
  \vspace{4pt}
  \begin{columns}[t,onlytextwidth]
    \begin{column}{0.32\textwidth}
      \textcolor{bleu303}{\bfseries\small Partie I}\\
      \textcolor{bleu303}{\small Le modèle}
      \vspace{6pt}
      \begin{itemize}\itemsep6pt\small
        \item[\textcolor{marron871}{$\bullet$}] Poisson : formule exacte
        \item[\textcolor{marron871}{$\bullet$}] Hawkes et la mémoire
        \item[\textcolor{marron871}{$\bullet$}] L'effet dominant
      \end{itemize}
    \end{column}
    \begin{column}{0.32\textwidth}
      \textcolor{violet2622}{\bfseries\small Partie II}\\
      \textcolor{violet2622}{\small Événements rares}
      \vspace{6pt}
      \begin{itemize}\itemsep6pt\small
        \item[\textcolor{marron871}{$\bullet$}] Monte-Carlo : impossible
        \item[\textcolor{marron871}{$\bullet$}] AMS : décomposer l'impossible
        \item[\textcolor{marron871}{$\bullet$}] La mémoire change tout
      \end{itemize}
    \end{column}
    \begin{column}{0.32\textwidth}
      \textcolor{vert575}{\bfseries\small Partie III}\\
      \textcolor{vert575}{\small Confrontation aux données}
      \vspace{6pt}
      \begin{itemize}\itemsep6pt\small
        \item[\textcolor{marron871}{$\bullet$}] Calibration MLE
        \item[\textcolor{marron871}{$\bullet$}] Bitcoin Black Thursday
      \end{itemize}
    \end{column}
  \end{columns}

  \vfill
  \begin{center}
    \textcolor{gris10}{\small\itshape
      Fil conducteur : comment mesurer la probabilité de l'impensable ?}
  \end{center}
\end{frame}

% =============================================================================
\partframe{bleu303}{Partie I}{Construire le modèle}
% =============================================================================

% ─────────────────────────────────────────────────────────────────────────────
\begin{frame}{Le carnet d'ordres comme problème de ruine}
% ─────────────────────────────────────────────────────────────────────────────
  \begin{columns}[c,onlytextwidth]
    \begin{column}{0.48\textwidth}
      \vspace{4pt}
      Un carnet électronique : deux files de quantités, bid et ask.
      Quand une file s'épuise, \textbf{le prix saute}.

      \bigskip
      On cherche deux quantités~:
      \[
        \Pb(\tau_a < \tau_b)
        \qquad \text{et} \qquad
        \E\bigl[\min(\tau_a,\tau_b)\bigr]
      \]

      \bigskip
      \textcolor{bleu303}{\bfseries Pourquoi c'est non-trivial~:}
      les files sont couplées, les arrivées ont de la mémoire,
      et les événements extrêmes sont rarissimes.
    \end{column}
    \begin{column}{0.48\textwidth}
      \centering
      \fig[width=0.92\linewidth]{trajectory_plane.png}
      \vspace{3pt}
      {\scriptsize\color{gris10}
        La trajectoire $(q_b,q_a)$. Le premier côté à zéro
        détermine la direction du saut de prix.}
    \end{column}
  \end{columns}
\end{frame}

% ─────────────────────────────────────────────────────────────────────────────
\begin{frame}{Résultat inattendu : la loi inverse-gaussienne est déjà opérationnelle}
% ─────────────────────────────────────────────────────────────────────────────
  \begin{columns}[c,onlytextwidth]
    \begin{column}{0.48\textwidth}
      \vspace{4pt}
      Dans la limite diffusive, la théorie prédit~:
      \[
        \tau \;\sim\; \IG\!\Bigl(\tfrac{q_0}{|\mu|},\;\tfrac{q_0^2}{\sigma^2}\Bigr)
      \]

      \bigskip
      \textcolor{bleu303}{\bfseries Ce qui surprend~:} cette approximation
      asymptotique est \emph{déjà} valide pour $q_0 = 10$
      — fini, discret.

      \medskip
      Test KS~: $p = 0{,}097$ — non-rejeté à $5\,\%$.

      \bigskip
      \textcolor{bleu303}{\bfseries Ce que ça débloque~:} calculer
      $\Pb(\tau_a < \tau_b)$ \emph{sans simulation}, par quadrature.
    \end{column}
    \begin{column}{0.48\textwidth}
      \centering
      \fig[width=0.92\linewidth]{ig_distribution.png}
      \vspace{3pt}
      {\scriptsize\color{gris10}
        Densité IG (rouge) vs histogramme simulé ($N=50\,000$, bleu).
        Accord sub-pourcent dès $q_0=10$.}
    \end{column}
  \end{columns}
\end{frame}

% ─────────────────────────────────────────────────────────────────────────────
\begin{frame}{Hawkes : la mémoire épaissit les deux queues simultanément}
% ─────────────────────────────────────────────────────────────────────────────
  \begin{columns}[c,onlytextwidth]
    \begin{column}{0.48\textwidth}
      \vspace{4pt}
      Sur un vrai carnet, les annulations se déclenchent mutuellement.
      Le processus de Hawkes capture ça~:
      \[
        \lambda^-(t) = \mu^- + \alpha\!\sum_{s<t}e^{-\beta(t-s)}
      \]

      \bigskip
      \textcolor{bleu303}{\bfseries L'effet contre-intuitif~:}
      l'auto-excitation produit une queue gauche épaisse
      (avalanches) \textbf{ET} une queue droite épaisse
      (périodes calmes). Distribution bien plus étalée que Poisson.

      \bigskip
      \textcolor{rouge485}{\bfseries Biais structurel~:}
      la troncature $\max(\lambda^-,0)$ brise le bilan de flux.
      Biais de $+5{,}9\,\%$ sur $\bar\lambda^-$ — systématique,
      non-éliminable.
    \end{column}
    \begin{column}{0.48\textwidth}
      \centering
      \fig[width=0.92\linewidth]{hawkes_stopping_dist.png}
      \vspace{3pt}
      {\scriptsize\color{gris10}
        Distribution de $\tau$~: Hawkes (bleu) vs Poisson de fond (tirets).}
    \end{column}
  \end{columns}
\end{frame}

% ─────────────────────────────────────────────────────────────────────────────
\begin{frame}{Le couplage croisé est l'effet dominant : facteur $\div$\,2}
% ─────────────────────────────────────────────────────────────────────────────
  \begin{columns}[c,onlytextwidth]
    \begin{column}{0.44\textwidth}
      \centering
      \fig[width=0.95\linewidth]{comparison.png}
      \vspace{3pt}
      {\scriptsize\color{gris10}
        Densités empiriques des quatre régimes.}
    \end{column}
    \begin{column}{0.52\textwidth}
      \vspace{4pt}
      \begin{center}
      \renewcommand{\arraystretch}{1.3}
      \small
      \begin{tabular}{@{}lc@{}}
        \toprule
        Modèle & $\E[\min(\tau_a,\tau_b)]$ \\
        \midrule
        Poisson indépendant   & $12{,}4$ \\
        Min de deux Poisson   & $8{,}6$ \\
        Hawkes individuel     & $8{,}8$ \\
        \textcolor{bleu303}{\bfseries Hawkes couplé}
          & \textcolor{bleu303}{\bfseries$5{,}9$} \\
        \bottomrule
      \end{tabular}
      \end{center}

      \bigskip
      Les effets se \textbf{multiplient}, pas s'additionnent~:
      combinatoire ($\div 1{,}45$) $\times$ auto-excitation $\times$ couplage
      $= \mathbf{\div\,2{,}1}$.

      \bigskip
      \textcolor{gris10}{\small\itshape
        $\Rightarrow$ Un carnet qui s'épuise deux fois plus vite~:
        le couplage n'est pas une correction, c'est la dynamique dominante.}
    \end{column}
  \end{columns}
\end{frame}

% =============================================================================
\partframe{violet2622}{Partie II}{Quantifier l'extrême rare}
% =============================================================================

% ─────────────────────────────────────────────────────────────────────────────
\begin{frame}{Flash crash : le Monte-Carlo naïf est inutilisable}
% ─────────────────────────────────────────────────────────────────────────────
  \begin{columns}[c,onlytextwidth]
    \begin{column}{0.50\textwidth}
      \vspace{4pt}
      Un flash crash de profondeur $k$ = $k$ bougies négatives consécutives.
      La probabilité décroît exponentiellement en $k$.

      \bigskip
      CV $= \sqrt{(1-p)/(Mp)}$ diverge en $1/\sqrt{p}$~:
      \begin{center}
      \small\renewcommand{\arraystretch}{1.25}
      \begin{tabular}{@{}lcc@{}}
        \toprule
        $k$ & $p$ & $M^*$ pour CV\,$=0{,}1$ \\
        \midrule
        $1$ & $0{,}14$ & $600$ \\
        $5$ & $10^{-7}$ & $10^9$ \\
        $8$ & $10^{-13}$ & $5\times10^{14}$ \\
        \bottomrule
      \end{tabular}
      \end{center}

      \vspace{6pt}
      À $10^9$ trajectoires/s, estimer $k=8$
      prendrait \textbf{15\,000 ans}.
    \end{column}
    \begin{column}{0.46\textwidth}
      \centering
      \vspace{8pt}
      \begin{tikzpicture}[scale=0.88]
        \draw[->,gray!50,thick] (0,0) -- (5.0,0)
          node[right,gray!60,font=\small] {$k$};
        \draw[->,gray!50,thick] (0,0) -- (0,4.2)
          node[above,gray!60,font=\small] {$-\log_{10}(p)$};
        \draw[violet2622,very thick,domain=0.3:4.5,samples=40]
          plot (\x,{0.88*\x+0.05});
        \fill[rouge485!12,rounded corners=3pt]
          (3.0,2.0) rectangle (4.8,4.1);
        \node[rouge485!80!black,font=\small,align=center]
          at (3.9,3.05) {inaccessible\\au MC};
        \foreach \k in {1,2,3,4}{
          \node[gray!60,font=\small,below] at (\k,0) {\k};
        }
        \node[gray!60,font=\scriptsize,left] at (0,0.9) {$10^{-1}$};
        \node[gray!60,font=\scriptsize,left] at (0,1.8) {$10^{-3}$};
        \node[gray!60,font=\scriptsize,left] at (0,2.7) {$10^{-7}$};
        \fill[violet2622] (1,0.93) circle (2.5pt);
        \fill[violet2622] (2,1.81) circle (2.5pt);
        \fill[rouge485!80!black] (3,2.69) circle (2.5pt);
        \fill[rouge485!80!black] (4,3.57) circle (2.5pt);
      \end{tikzpicture}
    \end{column}
  \end{columns}
\end{frame}

% ─────────────────────────────────────────────────────────────────────────────
\begin{frame}{AMS : décomposer l'impossible en produit de facile}
% ─────────────────────────────────────────────────────────────────────────────
  \begin{columns}[c,onlytextwidth]
    \begin{column}{0.50\textwidth}
      \vspace{4pt}
      \textbf{Idée~:} utiliser une fonction de score (profondeur atteinte)
      et décomposer l'événement en niveaux successifs~:
      \[
        \hat P_\mathrm{AMS}
        = \hat p_9 \times \hat p_8 \times \cdots \times \hat p_0
      \]
      Chaque $\hat p_k \approx 0{,}5$ — chaque étape est triviale.

      \bigskip
      \textcolor{bleu303}{\bfseries Point crucial~:} le ré-échantillonnage
      préserve l'état Hawkes complet $(\lambda_a^-,\lambda_b^-)$.
      La mémoire est transmise entre niveaux.

      \bigskip
      \textbf{Résultat~:} CV passe de $0{,}35$ $(N=50)$
      à $0{,}05$ $(N=3200)$.\\
      Atteint $p\sim 10^{-13}$ avec $N=2000$ particules.
    \end{column}
    \begin{column}{0.46\textwidth}
      \centering
      \fig[width=0.92\linewidth]{ams_convergence.png}
      \vspace{3pt}
      {\scriptsize\color{gris10}
        Pente log-log $\approx -0{,}5$~:
        convergence en $1/\sqrt{N}$ confirmée empiriquement.}
    \end{column}
  \end{columns}
\end{frame}

% ─────────────────────────────────────────────────────────────────────────────
\begin{frame}{La mémoire invalide le modèle i.i.d. par un facteur $10^4$}
% ─────────────────────────────────────────────────────────────────────────────
  \begin{columns}[c,onlytextwidth]
    \begin{column}{0.44\textwidth}
      \centering
      \fig[width=0.95\linewidth]{flash_crash_asym.png}
      \vspace{3pt}
      {\scriptsize\color{gris10}
        AMS Hawkes (bleu) vs Bernoulli i.i.d.\ (orange) — échelle log.}
    \end{column}
    \begin{column}{0.52\textwidth}
      \vspace{4pt}
      Si les cycles étaient indépendants~: $P(\text{FC}_k) = p_1^k$.
      \textbf{Avec la mémoire Hawkes, c'est faux.}

      \medskip
      Après un épuisement bid, l'intensité ask a été excitée.
      Le cycle suivant \textit{favorise} l'ask — les cycles ne sont pas i.i.d.

      \vspace{8pt}
      \begin{center}
      \small\renewcommand{\arraystretch}{1.3}
      \begin{tabular}{@{}crr@{}}
        \toprule
        $k$ & AMS Hawkes & Bernoulli $p_1^k$ \\
        \midrule
        2 & $6{,}1\times10^{-3}$ & $2{,}0\times10^{-2}$ \\
        5 & $1{,}1\times10^{-7}$ & $5{,}8\times10^{-6}$ \\
        8 & $\mathbf{1{,}9\times10^{-13}}$ & $1{,}7\times10^{-9}$ \\
        \bottomrule
      \end{tabular}
      \end{center}

      \vspace{6pt}
      \textcolor{rouge485}{\bfseries Facteur $>10^4$ à $k=8$.}
      Le modèle i.i.d.\ surestime massivement les crash profonds.
    \end{column}
  \end{columns}
\end{frame}

% ─────────────────────────────────────────────────────────────────────────────
\begin{frame}{Importance Sampling : validation croisée et borne de confiance}
% ─────────────────────────────────────────────────────────────────────────────
  \begin{columns}[c,onlytextwidth]
    \begin{column}{0.50\textwidth}
      \vspace{4pt}
      On tord les intensités~:
      $\tilde\lambda_a^- = e^{-\theta}\lambda_a^-$,
      $\tilde\lambda_b^- = e^\theta\lambda_b^-$,
      et on corrige par le rapport de vraisemblance (Girsanov).

      \bigskip
      Le balayage en $\theta$ révèle~:
      \begin{itemize}\itemsep4pt
        \item $\theta^* \approx 0{,}30$~: CV minimal, RE\,$=0{,}57$
        \item $\theta > 0{,}60$~: poids explosifs, estimateur dégénéré
      \end{itemize}

      \bigskip
      \textcolor{bleu303}{\bfseries Les deux méthodes donnent
      $\hat p \approx 0{,}139$}.\\
      L'accord AMS\,/\,IS valide le résultat~:
      ce n'est pas un artefact de méthode.
    \end{column}
    \begin{column}{0.46\textwidth}
      \centering
      \fig[width=0.92\linewidth]{is_comparison.png}
      \vspace{3pt}
      {\scriptsize\color{gris10}
        CV vs $\theta$~: minimum net à $\theta^*\approx0{,}30$,
        effondrement au-delà.}
    \end{column}
  \end{columns}
\end{frame}

% =============================================================================
\partframe{vert575}{Partie III}{Confronter le modèle aux données}
% =============================================================================

% ─────────────────────────────────────────────────────────────────────────────
\begin{frame}{Calibration MLE : ce qui est robuste et ce qui est fragile}
% ─────────────────────────────────────────────────────────────────────────────
  \begin{columns}[c,onlytextwidth]
    \begin{column}{0.50\textwidth}
      \vspace{4pt}
      Log-vraisemblance bivarée calculée en $O(n)$
      via les récurrences Hawkes.

      \bigskip
      \textcolor{bleu303}{\bfseries Robuste~:}
      les taux de fond $\mu_a^-$, $\mu_b^-$ sont retrouvés
      à $<0{,}5\,\%$ sur données synthétiques.

      \medskip
      \textcolor{rouge485}{\bfseries Fragile~:}
      $(\alpha,\beta)$ sont fortement corrélés sur séquences courtes.
      Des données tick-by-tick sont nécessaires pour les séparer.

      \vspace{8pt}
      \begin{center}
      \small\renewcommand{\arraystretch}{1.25}
      \begin{tabular}{@{}lcccc@{}}
        \toprule
        & $\mu_a^-$ & $\mu_b^-$ & $\alpha$ & $\beta$ \\
        \midrule
        Vrai & $3{,}0$ & $1{,}5$ & $0{,}35$ & $1{,}0$ \\
        MLE  & $3{,}01$ & $1{,}50$ & $0{,}37$ & $1{,}15$ \\
        \bottomrule
      \end{tabular}
      \end{center}
    \end{column}
    \begin{column}{0.46\textwidth}
      \centering
      \fig[width=0.92\linewidth]{mle_calibration.png}
      \vspace{3pt}
      {\scriptsize\color{gris10}
        Profil de vraisemblance~: un seul maximum bien identifié pour $\mu_a^-$.}
    \end{column}
  \end{columns}
\end{frame}

% ─────────────────────────────────────────────────────────────────────────────
\begin{frame}{Bitcoin Black Thursday : $p = 8{,}47\times10^{-6}$}
% ─────────────────────────────────────────────────────────────────────────────
  \begin{columns}[c,onlytextwidth]
    \begin{column}{0.50\textwidth}
      \vspace{4pt}
      12 mars 2020 : BTC perd \textbf{41\,\%} en 24h.
      5 bougies négatives consécutives observées.

      \bigskip
      On calibre sur $1\,696$ bougies horaires \emph{pré-crise}~:
      \[
        \hat\alpha \simeq 0
        \quad \Rightarrow \quad
        \text{régime quasi-Poisson avant la crise.}
      \]

      La crise est un \textbf{changement de régime},
      pas un régime stationnaire excité.

      \bigskip
      Probabilité AMS sous le régime pré-crise~:
      \[
        P(k=5) = \mathbf{8{,}47\times10^{-6}
        \;\approx\; \frac{1}{118\,000}}
      \]
      Événement attendu \emph{une fois tous les 13 ans} à ce régime.
    \end{column}
    \begin{column}{0.46\textwidth}
      \centering
      \fig[width=0.92\linewidth]{bitcoin_calibration.png}
      \vspace{3pt}
      {\scriptsize\color{gris10}
        Cours BTC, taux calibrés, et $P(k)$ AMS vs i.i.d.}
    \end{column}
  \end{columns}
\end{frame}

% =============================================================================
%  CONCLUSION
% =============================================================================

\begin{frame}{Ce qu'on retient}
  \begin{columns}[t,onlytextwidth]
    \begin{column}{0.46\textwidth}
      \vspace{4pt}
      \textcolor{bleu303}{\bfseries\small Trois résultats qui comptent}
      \begin{itemize}\itemsep8pt\small
        \item La loi IG fonctionne dès $q_0=10$~:
              formule analytique pour $\Pb(\tau_a<\tau_b)$
        \item Le couplage croisé divise $\E[\min(\tau)]$ par 2
        \item La mémoire Hawkes modifie $P(\text{flash crash})$
              d'un facteur $\mathbf{>10^4}$
      \end{itemize}

      \bigskip
      \textcolor{rouge485}{\bfseries\small Deux limites honnêtes}
      \begin{itemize}\itemsep6pt\small
        \item Biais de troncature de $\lambda^-$~: systématique,
              non-éliminable
        \item GameStop~: 70 obs.\ journalières insuffisantes
              pour calibrer $(\alpha,\beta)$
      \end{itemize}
    \end{column}
    \begin{column}{0.50\textwidth}
      \vspace{4pt}
      \textcolor{vert575}{\bfseries\small Perspectives}
      \begin{itemize}\itemsep6pt\small
        \item Données tick-by-tick (TARDIS, Kaiko)
              pour résoudre la corrélation $\alpha$-$\beta$
        \item Extension à $L>2$ niveaux~: architecture déjà en place
        \item Relier $\Pb(\tau_a<\tau_b)$ calibré
              à une surface de volatilité implicite
      \end{itemize}

      \vfill
      \begin{beamercolorbox}[sep=8pt]{block body}
        \centering\small
        La mémoire Hawkes n'est pas une correction marginale.\\[4pt]
        Elle change les probabilités de flash crash\\
        d'un facteur supérieur à $\mathbf{10^4}$.
      \end{beamercolorbox}
    \end{column}
  \end{columns}
\end{frame}

% =============================================================================
%  ANNEXES
% =============================================================================

{
\setbeamertemplate{footline}{}
\begin{frame}[plain,c,noframenumbering]
  \begin{tikzpicture}[overlay,remember picture]
    \fill[bleu303] (current page.south west) rectangle (current page.north east);
    \fill[marron871]
      (current page.north west) rectangle
      ([xshift=5pt]current page.south west);
    \node[white,font=\fontsize{28}{34}\selectfont\sffamily\bfseries]
      at (current page.center) {Annexes};
    \draw[marron871!60!white,line width=0.7pt]
      ([xshift=-2.5cm,yshift=-0.9cm]current page.center)
      -- ([xshift=2.5cm,yshift=-0.9cm]current page.center);
  \end{tikzpicture}
\end{frame}
}

% ─────────────────────────────────────────────────────────────────────────────
\begin{frame}{Annexe A — Surfaces 3D : Poisson vs Hawkes couplé}
% ─────────────────────────────────────────────────────────────────────────────
  \begin{columns}[c,onlytextwidth]
    \begin{column}{0.48\textwidth}
      \centering
      {\small Modèle de Poisson}\\[6pt]
      \fig[width=0.90\linewidth]{time_proba_3d_poisson.png}
    \end{column}
    \begin{column}{0.48\textwidth}
      \centering
      {\small Hawkes couplé}\\[6pt]
      \fig[width=0.90\linewidth]{hawkes_coupled_3d.png}
    \end{column}
  \end{columns}
  \vspace{6pt}
  \small\textcolor{gris10}{
    Le couplage croisé abaisse $\E[\min(\tau)]$ de $8{,}32$ à $5{,}83$
    au coin $(10,10)$ sans briser la symétrie diagonale.}
\end{frame}

% ─────────────────────────────────────────────────────────────────────────────
\begin{frame}{Annexe B — Comparaison AMS vs IS}
% ─────────────────────────────────────────────────────────────────────────────
  \begin{columns}[c,onlytextwidth]
    \begin{column}{0.50\textwidth}
      \vspace{4pt}
      $R=50$ réplications, $N=500$ trajectoires.

      \vspace{8pt}
      \begin{center}
      \small\renewcommand{\arraystretch}{1.3}
      \begin{tabular}{@{}lcc@{}}
        \toprule
        Estimateur & $\hat\sigma(\hat P)$ & RE \\
        \midrule
        AMS & $0{,}0171$ & $1{,}00$ \\
        IS ($\theta^*=0{,}30$) & $0{,}0129$ & $\mathbf{0{,}57}$ \\
        \bottomrule
      \end{tabular}
      \end{center}

      \vspace{8pt}
      IS plus précis à $N$ modéré.
      AMS reprend l'avantage pour $k\ge5$~:
      sa structure par niveaux est indispensable
      pour les probabilités $<10^{-7}$.
    \end{column}
    \begin{column}{0.46\textwidth}
      \centering
      \fig[width=0.92\linewidth]{clt_comparison.png}
      \vspace{3pt}
      {\scriptsize\color{gris10}
        Histogrammes sur $R=50$ répétitions~;
        largeur IC 95\,\% vs $N$.}
    \end{column}
  \end{columns}
\end{frame}

\end{document}
